% GENERATED FILE. Edit canonical lessons, then run npm run generate:books.
% canonical-source-hash: fnv1a64:341eaf4a8fb126ee
% canonical-lessons: FA-C253-piyade-shodan, FA-C253-vared-shodan, FA-C253-kharej-shodan, FA-C253-hel-dadan, FA-C253-andakhtan

\chapter{To get off, To enter, To go out, To push, To throw}
\label{ch:fa-to-get-off-to-enter-to-go-out-to-push-to-throw}
\hlchaptermodality{253}

\begin{chapteropening}
\textbf{By the end of this chapter:} \emph{I can say to get off, to enter, to go out, to push and to throw.}

Complete the last lesson of chapter 253: I can say to get off, to enter, to go out, to push and to throw.
\end{chapteropening}

\section[\texorpdfstring{\fa{پیاده} \fa{شدن} (piyâde shodan)}{پیاده شدن (piyâde shodan)}]{\fa{پیاده} \fa{شدن} (piyâde shodan) — to get off (a vehicle)}

\label{lesson:FA-C253-piyade-shodan}



\begin{quote}
Before the new one: say the Persian for a guard, then the Persian for a street sweeper.
\end{quote}

\subsection*{You'll want to know: \fa{پیاده} \fa{شدن}}
\textbf{\fa{پیاده} \fa{شدن}} — \emph{piyâde shodan} — \textquotedblleft{}to get off (a vehicle)\textquotedblright{}.

\begin{cousinweb}
\emph{piyâde shodan} is \textquotedblleft{}to become on foot\textquotedblright{}: \emph{piyâde} is inherited, from Middle Persian \emph{payādag}, \textquotedblleft{}foot soldier\textquotedblright{} (EtymDB, which also lists the sense \textquotedblleft{}chess pawn\textquotedblright{}). \emph{shodan}, \textquotedblleft{}to become\textquotedblright{}, is an inherited verb; it is usually traced to Middle Persian \emph{šudan}, \textquotedblleft{}to go\textquotedblright{}.
\end{cousinweb}

\begin{grammarlens}[title={What you've built}]
One more word for this chapter.
\end{grammarlens}

\subsection*{Your turn}
\begin{itemize}
\raggedright
  \item \emph{Say it:} \emph{piyâde shodan}
  \item \emph{Say it:} \emph{piyâde shodan}, once more
  \item \emph{Say it:} say \emph{piyâde shodan}
  \item \emph{Recall:} say the Persian for a guard, then the Persian for a street sweeper, then say \emph{piyâde shodan} again
\end{itemize}

\begin{tcolorbox}[breakable,colback=teal!4,colframe=teal!35!black,title={Before you move on}]
What does \fa{پیاده} \fa{شدن} mean? (\textquotedblleft{}To get off (a vehicle)\textquotedblright{}.) Say it once more.
\end{tcolorbox}
\section[\texorpdfstring{\fa{وارد} \fa{شدن} (vâred shodan)}{وارد شدن (vâred shodan)}]{\fa{وارد} \fa{شدن} (vâred shodan) — to enter}

\label{lesson:FA-C253-vared-shodan}



\begin{quote}
Before the new one: say the Persian for a street sweeper, then the Persian for to get off.
\end{quote}

\subsection*{You'll want to know: \fa{وارد} \fa{شدن}}
\textbf{\fa{وارد} \fa{شدن}} — \emph{vâred shodan} — \textquotedblleft{}to enter\textquotedblright{}.

\begin{cousinweb}
\emph{vâred shodan} pairs \emph{vâred}, borrowed from Arabic \emph{wārid}, \textquotedblleft{}arriving, coming in\textquotedblright{} (root \emph{w-r-d}, Lane), with the Persian verb \emph{shodan}, \textquotedblleft{}to become\textquotedblright{}. \emph{shodan}, \textquotedblleft{}to become\textquotedblright{}, is an inherited verb; it is usually traced to Middle Persian \emph{šudan}, \textquotedblleft{}to go\textquotedblright{}.
\end{cousinweb}

\begin{grammarlens}[title={What you've built}]
One more word for this chapter.
\end{grammarlens}

\subsection*{Your turn}
\begin{itemize}
\raggedright
  \item \emph{Say it:} \emph{vâred shodan}
  \item \emph{Say it:} \emph{vâred shodan}, once more
  \item \emph{Say it:} say \emph{vâred shodan}
  \item \emph{Recall:} say the Persian for a street sweeper, then the Persian for to get off, then say \emph{vâred shodan} again
\end{itemize}

\begin{tcolorbox}[breakable,colback=teal!4,colframe=teal!35!black,title={Before you move on}]
What does \fa{وارد} \fa{شدن} mean? (\textquotedblleft{}To enter\textquotedblright{}.) Say it once more.
\end{tcolorbox}
\section[\texorpdfstring{\fa{خارج} \fa{شدن} (khârej shodan)}{خارج شدن (khârej shodan)}]{\fa{خارج} \fa{شدن} (khârej shodan) — to go out, to exit}

\label{lesson:FA-C253-kharej-shodan}



\begin{quote}
Before the new one: say the Persian for to get off, then the Persian for to enter.
\end{quote}

\subsection*{You'll want to know: \fa{خارج} \fa{شدن}}
\textbf{\fa{خارج} \fa{شدن}} — \emph{khârej shodan} — \textquotedblleft{}to go out, to exit\textquotedblright{}.

\begin{cousinweb}
\emph{khârej shodan} pairs \emph{khârej}, borrowed from Arabic \emph{ḫārij}, \textquotedblleft{}outside, going out\textquotedblright{} (root \emph{ḫ-r-j}, Lane), with the Persian verb \emph{shodan}, \textquotedblleft{}to become\textquotedblright{}. \emph{shodan}, \textquotedblleft{}to become\textquotedblright{}, is an inherited verb; it is usually traced to Middle Persian \emph{šudan}, \textquotedblleft{}to go\textquotedblright{}.
\end{cousinweb}

\begin{grammarlens}[title={What you've built}]
One more word for this chapter.
\end{grammarlens}

\subsection*{Your turn}
\begin{itemize}
\raggedright
  \item \emph{Say it:} \emph{khârej shodan}
  \item \emph{Say it:} \emph{khârej shodan}, once more
  \item \emph{Say it:} say \emph{khârej shodan}
  \item \emph{Recall:} say the Persian for to get off, then the Persian for to enter, then say \emph{khârej shodan} again
\end{itemize}

\begin{tcolorbox}[breakable,colback=teal!4,colframe=teal!35!black,title={Before you move on}]
What does \fa{خارج} \fa{شدن} mean? (\textquotedblleft{}To go out, to exit\textquotedblright{}.) Say it once more.
\end{tcolorbox}
\section[\texorpdfstring{\fa{هل} \fa{دادن} (hel dâdan)}{هل دادن (hel dâdan)}]{\fa{هل} \fa{دادن} (hel dâdan) — to push}

\label{lesson:FA-C253-hel-dadan}



\begin{quote}
Before the new one: say the Persian for to enter, then the Persian for to go out.
\end{quote}

\subsection*{You'll want to know: \fa{هل} \fa{دادن}}
\textbf{\fa{هل} \fa{دادن}} — \emph{hel dâdan} — \textquotedblleft{}to push\textquotedblright{}.

\begin{cousinweb}
\emph{hel dâdan} pairs \emph{hel}, \textquotedblleft{}a push\textquotedblright{}, with \emph{dâdan}, \textquotedblleft{}to give\textquotedblright{}. Where \emph{hel} comes from is not settled. \emph{dâdan}, \textquotedblleft{}to give\textquotedblright{}, is an inherited verb; EtymDB takes it back to the Indo-European root \emph{*deh\textsubscript{3}-}, \textquotedblleft{}to give\textquotedblright{}.
\end{cousinweb}

\begin{grammarlens}[title={What you've built}]
One more word for this chapter.
\end{grammarlens}

\subsection*{Your turn}
\begin{itemize}
\raggedright
  \item \emph{Say it:} \emph{hel dâdan}
  \item \emph{Say it:} \emph{hel dâdan}, once more
  \item \emph{Say it:} say \emph{hel dâdan}
  \item \emph{Recall:} say the Persian for to enter, then the Persian for to go out, then say \emph{hel dâdan} again
\end{itemize}

\begin{tcolorbox}[breakable,colback=teal!4,colframe=teal!35!black,title={Before you move on}]
What does \fa{هل} \fa{دادن} mean? (\textquotedblleft{}To push\textquotedblright{}.) Say it once more.
\end{tcolorbox}
\section[\texorpdfstring{\fa{انداختن} (andâkhtan)}{انداختن (andâkhtan)}]{\fa{انداختن} (andâkhtan) — to throw}

\label{lesson:FA-C253-andakhtan}



\begin{quote}
Before the new one: say the Persian for to go out, then the Persian for to push.
\end{quote}

\subsection*{You'll want to know: \fa{انداختن}}
\textbf{\fa{انداختن}} — \emph{andâkhtan} — \textquotedblleft{}to throw\textquotedblright{}.

\begin{cousinweb}
\emph{andâkhtan} is a native Persian word, inherited rather than borrowed; it is usually traced to Middle Persian \emph{handāxtan}, \textquotedblleft{}to throw; to measure\textquotedblright{}.
\end{cousinweb}

\begin{grammarlens}[title={What you've built}]
One more word for this chapter.
\end{grammarlens}

\subsection*{Your turn}
\begin{itemize}
\raggedright
  \item \emph{Say it:} \emph{andâkhtan}
  \item \emph{Say it:} \emph{andâkhtan}, once more
  \item \emph{Say it:} say \emph{andâkhtan}
  \item \emph{Recall:} say the Persian for to go out, then the Persian for to push, then say \emph{andâkhtan} again
\end{itemize}

\begin{tcolorbox}[breakable,colback=teal!4,colframe=teal!35!black,title={Before you move on}]
What does \fa{انداختن} mean? (\textquotedblleft{}To throw\textquotedblright{}.) Say it once more.
\end{tcolorbox}
